% ============================================================
% Slide 13-16 (bản 45 phút) -- rút gọn từ part3.tex, part4.tex
% ============================================================

% --- Slide 13 (part3.tex: Gradient tập trung tại biên Gaussian) ---
\begin{frame}{Gradient tập trung tại biên Gaussian}
    \begin{columns}[T]
        \begin{column}{0.5\textwidth}
            \begin{itemize}
                \item Gradient view-space lớn nhất xuất hiện ở vùng \textbf{biên/rìa} của Gaussian -- nơi ảnh render lệch nhiều nhất so với ảnh thật.
                \item Gradient tại biên chính là \textbf{tín hiệu} quyết định \textbf{densify} (thêm Gaussian mới) ở giai đoạn sau.
                \item Vùng gradient thấp $\Rightarrow$ Gaussian đã khớp tốt, không cần chia tách thêm.
            \end{itemize}
        \end{column}
        \begin{column}{0.5\textwidth}
            \centering
            \includegraphics[width=0.95\linewidth]{ch6_edge_gradient.png}
            \captionof{figure}{Gradient tập trung ở biên Gaussian}
        \end{column}
    \end{columns}
\end{frame}

% --- Slide 14 (part3.tex: Tổng kết vòng lặp huấn luyện) ---
\begin{frame}{Tổng kết vòng lặp huấn luyện}
    \begin{center}
        \footnotesize
        \begin{tikzpicture}[
            node distance=1.1cm and 1.6cm,
            box/.style={draw, rounded corners, fill=blue!8, align=center, minimum width=2.6cm, minimum height=1cm},
            ->, thick
        ]
            \node[box] (render) {1. Render};
            \node[box, right=of render] (loss) {2. Tính loss\\ (L1 + D-SSIM)};
            \node[box, right=of loss] (backward) {3. Backward\\ (lan gradient)};
            \node[box, below=of loss] (step) {4. optimizer.step()\\ (sparse Adam)};
            \node[box, left=of step] (dp) {5. Densify /\\ Prune (định kỳ)};

            \draw (render) -- (loss);
            \draw (loss) -- (backward);
            \draw (backward) |- (step);
            \draw (step) -- (dp);
            \draw (dp) |- (render);
        \end{tikzpicture}
    \end{center}
    \begin{itemize}
        \item Một iteration: \textbf{render} $\to$ \textbf{loss} $\to$ \textbf{backward} $\to$ \textbf{optimizer.step()} $\to$ (định kỳ) \textbf{densify/prune} $\to$ lặp lại đến hội tụ.
    \end{itemize}
\end{frame}

% --- Slide 15 (part4.tex: Tổng quan Adaptive Density Control (ADC)) ---
\begin{frame}{Tổng quan Adaptive Density Control (ADC)}
    \begin{itemize}
        \item Point cloud khởi tạo từ COLMAP thường \textbf{quá thưa} ở vùng thiếu chi tiết, đồng thời có Gaussian \textbf{quá to} che khuất chi tiết nhỏ.
        \item ADC chạy định kỳ mỗi \texttt{densification\_interval} iteration để điều chỉnh mật độ Gaussian:
        \begin{itemize}
            \item \textbf{Densify}: nhân bản (clone) hoặc tách (split) Gaussian thiếu hình học.
            \item \textbf{Prune}: cắt tỉa Gaussian thừa, mờ, vô dụng.
        \end{itemize}
        \item Mục tiêu: cân bằng chất lượng tái tạo và số lượng Gaussian (bộ nhớ, tốc độ render).
    \end{itemize}
    \vspace{0.2cm}
    \centering
    \begin{tikzpicture}[node distance=1.4cm, every node/.style={font=\footnotesize}]
        \node[draw, rounded corners, fill=blue!10] (adc) {ADC};
        \node[draw, rounded corners, fill=green!10, below left=0.9cm and 1.6cm of adc] (clone) {Clone};
        \node[draw, rounded corners, fill=orange!10, below=1.2cm of adc] (split) {Split};
        \node[draw, rounded corners, fill=red!10, below right=0.9cm and 1.6cm of adc] (prune) {Prune};
        \draw[->] (adc) -- (clone);
        \draw[->] (adc) -- (split);
        \draw[->] (adc) -- (prune);
    \end{tikzpicture}
\end{frame}

% --- Slide 16 (part4.tex: Tích luỹ gradient view-space) ---
\begin{frame}{Tích luỹ gradient view-space}
\begin{columns}[T]
\begin{column}{0.52\textwidth}
\footnotesize
    \begin{itemize}
        \item Với mỗi Gaussian, gradient vị trí 2D (view-space) được \textbf{cộng dồn} qua nhiều lần render từ các camera khác nhau.
        \item Đồng thời đếm \texttt{denom}: số lần Gaussian đó thực sự visible trong batch.
        \item Gradient trung bình dùng làm tín hiệu densify:
        \[
            \bar{g} = \frac{\sum_{v} \nabla_{\text{view}} }{\text{denom}}
        \]
        \item $\bar g$ lớn $\Rightarrow$ ứng viên clone/split.
    \end{itemize}
\end{column}
\begin{column}{0.48\textwidth}
    \centering
    \includegraphics[width=0.95\linewidth]{fig_02_accum_denom.png}
\end{column}
\end{columns}
\end{frame}
